\documentclass[12pt,a4paper]{ctexart}

% ── 编码与字体 ──
\usepackage{amsmath,amssymb}

% ── 页面设置 ──
\usepackage[margin=2.5cm]{geometry}
\usepackage{setspace}
\onehalfspacing

% ── 插图与表格 ──
\usepackage{graphicx}
\usepackage{float}
\usepackage{caption}
\usepackage{booktabs}
\usepackage{array}
\usepackage{longtable}

% ── 颜色与链接 ──
\usepackage[colorlinks=true,linkcolor=blue,urlcolor=blue]{hyperref}

% ── 代码高亮 ──
\usepackage{listings}
\usepackage{xcolor}

\lstset{
    language=C++,
    basicstyle=\small\ttfamily,
    backgroundcolor=\color{gray!5},
    frame=single,
    breaklines=true,
    numbers=left,
    numberstyle=\tiny,
    keywordstyle=\color{blue},
    commentstyle=\color{green!50!black},
    stringstyle=\color{red!50!orange},
    tabsize=4,
    captionpos=b
}

\lstdefinestyle{csharp}{
    language=[Sharp]C,
    basicstyle=\small\ttfamily,
    backgroundcolor=\color{gray!5},
    frame=single,
    breaklines=true,
    numbers=left,
    numberstyle=\tiny,
    keywordstyle=\color{blue},
    commentstyle=\color{green!50!black},
    stringstyle=\color{red!50!orange},
    tabsize=4,
    captionpos=b
}

\lstdefinestyle{json}{
    language=TeX,
    basicstyle=\small\ttfamily,
    backgroundcolor=\color{gray!5},
    frame=single,
    breaklines=true,
    numbers=none,
    tabsize=2,
    captionpos=b
}

\title{\vspace{-2cm}符玄桌面宠物 --- 技术文档}
\author{符玄桌宠项目}
\date{V2.1 · 2026-07-09}

\begin{document}

\maketitle
\thispagestyle{empty}

\begin{abstract}
本文档为符玄桌面宠物项目的完整技术报告（V2.1）。项目基于 Tuanjie 2022.3.62t7（团结引擎）与 Live2D Cubism SDK 5-r.4，
构建了一个运行于 Windows 桌面的 AI 驱动虚拟角色。系统深度集成 DeepSeek Chat API（含 Function Calling 37+ 工具）、
GLM-4V 多模态视觉模型、以及多层级感知与记忆系统，实现了从自然语言交互到具身动作表达的完整闭环。
V2.1 引入了全新的悬浮球（BallPanel）与右侧 Widgets 面板（RightPanel）交互体系，取代了旧版右键菜单，
涵盖物理模拟、Live2D 参数化动画、DWM 透明窗口、AI 对话、天气感知、便签提醒等核心功能。
\end{abstract}

\tableofcontents
\newpage

% ═══════════════════════════════════════════
\section{项目概述}
% ═══════════════════════════════════════════

\subsection{项目背景}

原桌面宠物程序基于 Windows GDI+ 分层窗口技术，使用静态 PNG 图片硬切换模拟动画。
本项目将其迁移至 Unity 引擎，利用 Live2D Cubism SDK 实现参数化变形动画，
获得更流畅的视觉效果和更丰富的交互能力，并引入 AI 对话、环境感知、具身动作生成等高级特性。
版本 V2.1 对交互界面进行了彻底重构，以悬浮球辐射菜单和右侧 Widgets 面板替代了传统右键菜单。

\subsection{技术栈}

\begin{itemize}
    \item \textbf{引擎}：Tuanjie 2022.3.62t7（团结引擎）
    \item \textbf{角色渲染}：Live2D Cubism SDK 5-r.4（参数化变形 + 实时物理模拟）
    \item \textbf{3D 渲染}：Unity Animator + FBX（预留）
    \item \textbf{AI 对话}：DeepSeek Chat API + Function Calling（37+ 法术工具）
    \item \textbf{视觉模型}：GLM-4V（智谱）--- 截图分析 + 动作视觉验证
    \item \textbf{天气数据}：QWeather API（优先） / wttr.in（回退，双源）
    \item \textbf{文件搜索}：Everything CLI (es.exe) / 递归回退
    \item \textbf{推送通知}：Server酱³
    \item \textbf{窗口管理}：Win32 API（DWM / Shell\_NotifyIcon / UIAutomation）
    \item \textbf{构建脚本}：PowerShell + Unity Batchmode
    \item \textbf{平台}：Windows 10/11 64位
\end{itemize}

\subsection{核心指标}

\begin{itemize}
    \item 渲染帧率：20--60\,fps（自适应三档降频）
    \item AI 响应延迟：1--5\,s（取决于网络与工具调用轮次）
    \item 内存管理：800\,MB 触发 GC，1200\,MB 强制 GC
    \item 系统内存监控：85\% 预警 GC，93\% 紧急降质
    \item 工具调用：37+ 注册工具，最多 5 轮回环
    \item 空闲动作：12 种 JSON 配置驱动 + 2 种硬编码特效
    \item 动作验证：19 个测试用例，GLM-4V 视觉评分 + 闭环学习
    \item 长期记忆：核心事实 + Top-5 重要 + 近期琐事，JSON 持久化
\end{itemize}

% ═══════════════════════════════════════════
\section{系统架构}
% ═══════════════════════════════════════════

\subsection{分层架构}

系统采用分层架构设计，从底层到上层依次为：

\begin{lstlisting}[frame=single,numbers=none,style=csharp]
+----------------------------------------------------------+
|                     UI 表示层                              |
|   BallPanel(悬浮) | RightPanel(右侧) | ChatBubble          |
|   BottomInputBar | DebugWindow                             |
+----------------------------------------------------------+
|                     AI 核心层                              |
|   ChatManager | ToolCallInvoker | IdleChatGenerator       |
|   MotionTranslator | MotionAgent | MotionVerifier         |
+----------------------------------------------------------+
|                     感知层                                 |
|   ActivityTracker | BrowserTabReader | PetMemory          |
|   TimeWeatherController | PerformanceMonitor             |
+----------------------------------------------------------+
|                     物理层                                 |
|   DesktopPet (重力/碰撞/地面检测) | DragHandler           |
+----------------------------------------------------------+
|                     渲染层                                 |
|   HybridRenderer -> Live2DRenderer / Model3DRenderer     |
+----------------------------------------------------------+
|                     系统层                                 |
|   WindowOverlay (DWM) | SystemTrayManager | ServerPoll   |
+----------------------------------------------------------+
\end{lstlisting}

\subsection{执行顺序}

Unity 脚本执行顺序通过 \texttt{DefaultExecutionOrder} 属性严格控制：

\begin{enumerate}
    \item \texttt{DesktopPet.Update()} (order=0) --- 物理更新、状态转换、行走相位
    \item \texttt{CubismPhysicsController.LateUpdate()} (order=800) --- 衣服/头发/裙子物理模拟
    \item \texttt{Live2DRenderer.LateUpdate()} (order=801) --- 覆盖被物理重置的参数、空闲动画、交互反馈
\end{enumerate}

顺序 801 确保所有 Live2D 参数在 Cubism Physics 运算后写入，避免物理覆盖关键参数（如手臂摆幅）。

\subsection{依赖注入模式}

组件间通过 Unity Inspector 序列化引用或单例模式连接。全局单例包括：
\texttt{ChatConfig}（环境变量静态访问）、\texttt{PetMemory.Instance}（忆境）、
\texttt{PetConfig.Instance}（天机簿）。

% ═══════════════════════════════════════════
\section{渲染管道}
% ═══════════════════════════════════════════

\subsection{Live2D Cubism 渲染}

\subsubsection{模型加载双保险}

\texttt{Live2DRenderer} 使用两级加载策略：第一级加载 \texttt{AssetDatabase.LoadAssetAtPath}（Editor 模式），
失败时第二级通过 \texttt{Resources.Load} 降级（构建模式）。

\subsubsection{参数映射系统}

\texttt{Live2DParameterMapper} 建立语义参数名到 Cubism 参数 ID 的双向映射，涵盖约 80+ 个参数：

\begin{lstlisting}[style=csharp,caption={参数映射示例},label={lst:param-mapping}]
mapper.RegisterParameter("ParamAngleX", "head_angle_x");
mapper.RegisterParameter("ParamBodyAngleX", "body_angle_x");
mapper.SetParameterValue("head_angle_x", 15f);  // 头左倾 15 度
\end{lstlisting}

参数按身体部位分组：头部（3 轴旋转）、身体（3 轴倾角）、眼睛（开合/微笑/眼球）、
眉毛（上抬/集中）、嘴巴（形态/张合）、手臂（左右多段）、呼吸。

\subsubsection{Perlin 噪声驱动空闲微动}

空闲状态由 Perlin 噪声驱动 7 组微动参数，使用不同的种子值确保互不关联：

\[
\text{breath} = (\text{PerlinNoise}(t, 0) - 0.5) \times \text{Amplitude}
\]

\begin{table}[H]
\centering
\caption{Perlin 噪声微动参数}
\label{tab:perlin}
\begin{tabular}{lcc}
\toprule
\textbf{参数} & \textbf{通道 (seed)} & \textbf{振幅} \\
\midrule
ParamBreath & (t, 0) & BREATH\_AMPLITUDE \\
ParamBodyAngleX & (t, 1) & BODY\_SWAY\_X \\
ParamBodyAngleY & (t, 2) & BODY\_SWAY\_Y \\
ParamBodyAngleZ & (t, 3) & BODY\_SWAY\_Z \\
ParamAngleX & (t, 4) & HEAD\_X \\
ParamAngleY & (t, 5) & HEAD\_Y \\
ParamEyeBallX & (t+offset, 6) & EYE\_X \\
ParamEyeBallY & (t+offset, 7) & EYE\_Y \\
\bottomrule
\end{tabular}
\end{table}

\subsubsection{天气表情联动}

天气状态自动影响空闲表情基调：

\begin{table}[H]
\centering
\caption{天气表情联动}
\label{tab:weather-face}
\begin{tabular}{ll}
\toprule
\textbf{天气} & \textbf{表情效果} \\
\midrule
晴/多云 & ParamMouthForm 微笑 +0.2 \\
雨/阴 & ParamBrowRY/LY 眉抬 +4，MouthForm 微嘟 +0.2 \\
雷暴 & 同上 + ParamEyeLOpen 微睁 \\
雪 & ParamMouthOpenY 微张 +0.4，EyeLOpen 微睁 +0.2 \\
夜晚 (18--22点) & EyeLOpen 垂 +0.07 \\
深夜 (22--5点) & EyeLOpen 垂 +0.15 \\
\bottomrule
\end{tabular}
\end{table}

\subsection{DWM 透明窗口}

窗口透明使用 Windows DWM API 实现：

\begin{lstlisting}[style=csharp,caption={DWM 玻璃层扩展},label={lst:dwm}]
// 1. 设置窗口样式
SetWindowLong(hwnd, GWL_EXSTYLE, WS_EX_LAYERED | WS_EX_TRANSPARENT);

// 2. 扩展玻璃框架（黑色区域 = 透明）
MARGINS margins = new MARGINS { cxLeftWidth = -1 };
DwmExtendFrameIntoClientArea(hwnd, ref margins);

// 3. 设置背景纯黑透明
Camera.backgroundColor = new Color(0, 0, 0, 0);
\end{lstlisting}

关键特性：
\begin{itemize}
    \item \textbf{无绿边}：DWM 玻璃层与纯黑背景配合，消除色键抠图的半透明边缘伪影
    \item \textbf{分层渲染}：Live2D 渲染到 RenderTexture，通过 RawImage 显示在 Layer 31
    \item \textbf{点击穿透}：每帧动态切换 \texttt{WS\_EX\_TRANSPARENT}
    \item \textbf{DWM 崩溃保护}：连续 5 次重建失败则跳过 \texttt{DwmExtendFrameIntoClientArea}
\end{itemize}

% ═══════════════════════════════════════════
\section{AI 对话系统}
% ═══════════════════════════════════════════

\subsection{DeepSeek Chat API 集成}

\subsubsection{系统提示词注入}

\texttt{ChatManager.BuildSystemPrompt()} 在每次对话前动态构建 System Prompt，
注入内容包括：当前时间、法眼活动追踪摘要、浏览器标签页列表、
长期记忆格式化文本、Live2D 参数语义知识、具身动作生成能力描述。

\subsubsection{Function Calling 架构}

37+ 个注册工具通过 JSON Schema 暴露给 DeepSeek，覆盖网页搜索、截图分析、
系统控制、文件搜索（Everything 毫秒级）、便签管理、课表查询、动作/表情播放、
长期记忆操作、应用/文件夹启动、用户状态查询等。

\subsubsection{多轮工具调用循环}

最多 5 轮的工具调用回环：

\begin{lstlisting}[style=csharp,caption={多轮工具调用循环},label={lst:tool-loop}]
SendMessage()
  └── SendRequestCoroutine()
      ├── 1. 发送请求（含 system prompt + 历史）
      ├── 2. 解析 response
      ├── 3. 如果有 tool_calls:
      │     ├── 检查 round < 5
      │     ├── 遍历 tool_calls:
      │     │     ├── IsCoroutineTool()? -> 协程异步执行
      │     │     └── 同步工具 -> Execute()
      │     ├── 收集结果，回到 1（下一轮）
      └── 4. 无 tool_calls -> 处理回复文本
\end{lstlisting}

\subsubsection{异步工具执行}

5 个可能耗时的工具通过协程异步执行，每帧检查 \texttt{task.IsCompleted}，不阻塞主线程：

\begin{itemize}
    \item \texttt{query\_exams} / \texttt{query\_scores} / \texttt{query\_schedule} / \texttt{query\_user\_status}
          --- HTTP GET 请求课表小程序服务端
    \item \texttt{search\_files} --- Everything CLI 全盘搜索（毫秒级）或递归回退
\end{itemize}

\subsection{工具索引}

\begin{longtable}{lll}
\caption{工具列表（37+）}
\label{tab:tools}\\
\toprule
\textbf{工具} & \textbf{功能} & \textbf{类型} \\
\midrule
\endfirsthead
\multicolumn{3}{c}{\tablename\ \thetable\ 续前页} \\
\toprule
\textbf{工具} & \textbf{功能} & \textbf{类型} \\
\midrule
\endhead
\bottomrule
\endfoot
get\_time & 获取当前时间 & 同步 \\
open\_url & 打开网页/文件/文件夹 & 同步 \\
open\_app & 通过 PATH/StartMenu 启动应用 & 同步 \\
search\_web & 联网搜索（Bing） & 同步 \\
get\_weather & 查询天气 & 同步 \\
take\_screenshot & 截屏 + GLM-4V 分析 & 协程异步 \\
set\_volume / mute & 音量控制 & 同步 \\
get\_system\_info & 获取系统信息 & 同步 \\
lock\_screen & 锁屏 & 同步 \\
system\_power & 关机/重启/睡眠 & 同步（需确认） \\
run\_command & 执行 CMD & 同步（需确认） \\
get\_clipboard / set\_clipboard & 剪贴板读写 & 同步 \\
get\_memories / write\_memory & 长期记忆操作 & 同步 \\
start\_conversation & 发起聊天 & 同步 \\
messenger & 模拟回复 & 同步 \\
write\_note & 写便签 & 同步 \\
send\_notification & 桌面通知 & 同步 \\
notify & 推送通知（Server酱³） & 同步 \\
get\_pet\_status & 获取宠物状态 & 同步 \\
get\_mouse\_pos & 获取鼠标坐标 & 同步 \\
list\_files & 列出目录文件 & 同步 \\
set\_expression & 切换面部表情 & 同步 \\
play\_action / stop\_action & 播放/停止动作 & 同步 \\
open\_folder & 打开资源管理器 & 同步 \\
query\_exams & 查询考试安排 & 协程异步 \\
query\_scores & 查询成绩 & 协程异步 \\
query\_schedule & 查询课表 & 协程异步 \\
query\_user\_status & 查询学业概览 & 协程异步 \\
search\_files & 全盘文件搜索（Everything） & 协程异步 \\
inspect\_motion\_memory & 查看动作记忆 & 同步 \\
generate\_motion & LLM 生成自定义动作 & 协程异步 \\
get\_system\_status & 获取系统状态 & 同步 \\
show\_reminder / delete\_reminder & 便签管理 & 同步 \\
\bottomrule
\end{longtable}

\subsection{句子队列与优先级气泡}

长回复被拆分为单个句子，以打字机风格逐句显示。\texttt{ChatBubble} 实现了三级优先级抢占机制：

\begin{table}[H]
\centering
\caption{气泡优先级系统}
\label{tab:priority}
\begin{tabular}{lll}
\toprule
\textbf{优先级} & \textbf{用途} & \textbf{覆盖规则} \\
\midrule
Low & 闲话、定时问候、待机气泡 & 可被任何高优消息覆盖 \\
Normal & 提醒、交互回应 & 不可被 Low 覆盖 \\
High & AI 主动回复 & 不可被任何消息覆盖 \\
\bottomrule
\end{tabular}
\end{table}

\subsection{自动闲聊系统}

\texttt{IdleChatGenerator} 在无交互 60--120\,s 后触发，通过 DeepSeek API 批量预生成
（每次 10 条），本地 JSON 缓存。按时间周期自动刷新缓存，API 不可用时回退硬编码问候语。

% ═══════════════════════════════════════════
\section{具身动作系统}
% ═══════════════════════════════════════════

\subsection{架构演进}

动作系统经过 10 个阶段的持续演化：

\begin{itemize}
    \item \textbf{Phase 1--5}：硬编码 10 个空闲动作（直接写在 Live2DRenderer 中）
    \item \textbf{Phase 6}：动作参数提取为 KNOWN\_PATTERNS 常量
    \item \textbf{Phase 7}：7 个动作迁移到 JSON 配置 + \texttt{IdleActionScheduler}
    \item \textbf{Phase 8}：\texttt{MotionTranslator} --- LLM 翻译自然语言到关键帧
    \item \textbf{Phase 9 (N32d)}：闭环自评 --- GLM-4V 视觉验证 + \texttt{MotionMemoryManager} 优化
    \item \textbf{Phase 10 (V2.1)}：动作系统重构到 \texttt{Live2DFramework/ActionAgent/} 子目录，
          \texttt{MotionAgent} 引入自主动作决策（4 级密度），\texttt{IdleActionScheduler} 增加时段/天气权重调整，
          支持 12 种空闲动作
\end{itemize}

\subsection{空闲动作系统}

\subsubsection{JSON 配置驱动}

空闲动作定义在 \texttt{idle\_actions.json} 中，每个动作包含多阶段参数目标和插值曲线：

\begin{lstlisting}[style=json,caption={空闲动作 JSON 配置},label={lst:idle-json}]
{
  "formatVersion": "v2",
  "actions": [{
    "id": 1,
    "name": "head_tilt",
    "displayName": "歪头",
    "weight": 10,
    "cooldown": 8.0,
    "phases": [
      { "duration": 0.5, "curve": "smooth",
        "targets": { "ParamAngleZ": 5 } },
      { "duration": 0.5, "curve": "smooth",
        "targets": { "ParamAngleZ": 0 } }
    ]
  }]
}
\end{lstlisting}

\subsubsection{动作列表}

\begin{table}[H]
\centering
\caption{空闲动作列表}
\label{tab:idle-actions}
\begin{tabular}{cllc}
\toprule
\textbf{ID} & \textbf{名称} & \textbf{权重} & \textbf{类型} \\
\midrule
1 & 歪头 & 10 & JSON \\
2 & 微笑 & 10 & JSON \\
3 & 挑眉 & 8 & JSON \\
4 & 星辉缠绕 & 5 & 硬编码 \\
5 & 伸懒腰 & 8 & JSON \\
6 & 委屈 & 6 & JSON \\
7 & 法阵展开 & 4 & 硬编码 \\
8 & 害羞 & 6 & JSON \\
9 & 困惑 & 0 & JSON（AI 仅用）\\
10 & 爱心 & 5 & JSON \\
11 & 哭 & 4 & JSON \\
12 & 心跳 & 5 & JSON \\
\bottomrule
\end{tabular}
\end{table}

\subsubsection{插值曲线}

\texttt{MotionGenerator} 支持 6 种插值曲线：

\[
\begin{array}{ll}
\text{linear:} & f(t) = t \\
\text{smooth:} & f(t) = 3t^{2} - 2t^{3} \\
\text{easeOut:} & f(t) = 1 - (1-t)^{2} \\
\text{easeIn:} & f(t) = t^{2} \\
\text{hold:} & f(t) = 0 \\
\text{bounce:} & \text{分段弹性曲线}
\end{array}
\]

\subsection{LLM 动作翻译 (MotionTranslator)}

\subsubsection{核心流程}

\begin{enumerate}
    \item 构建 Body Schema --- 所有可用参数名、范围、语义描述
    \item 注入运动记忆 --- \texttt{MotionMemoryManager} 提供历史成功模板
    \item 调用 DeepSeek API（\texttt{temperature=0.3}，超时 30\,s）
    \item 解析 JSON 响应 → \texttt{MotionPlan}（关键帧序列）
    \item \texttt{MotionGenerator} 协程逐帧插值播放
\end{enumerate}

\subsubsection{SPECIAL PATTERNS}

内嵌在 System Prompt 中的姿势模板，帮助 DeepSeek 准确生成常见姿势：

\begin{table}[H]
\centering
\caption{SPECIAL PATTERNS}
\label{tab:special-patterns}
\begin{tabular}{ll}
\toprule
\textbf{模式} & \textbf{关键参数} \\
\midrule
叉腰 & arm\_right\_upper=0.7~0.8, arm\_right\_lower=0.15~0.25 \\
行礼 & body\_angle\_x=15~25, arm\_right\_lower=0.5~0.7 \\
歪头思考 & head\_angle\_z=15~25, eye\_ball\_y=-0.5~-0.7 \\
捂脸 & arm\_right\_upper=0.8~1.0, hand\_near\_face=1.0 \\
惊讶 & head\_angle\_y=8~15, eye\_l\_open=0.9~1.0 \\
缩团 & head\_angle\_y=-15~-25, arm 夹紧 \\
\bottomrule
\end{tabular}
\end{table}

\subsection{动作锁定机制}

AI 生成动作播放期间，通过 \texttt{\_aiControlLocked} 和 \texttt{\_actionLocked} 双锁防止被空闲动作或走路覆盖：
前者阻止 Perlin 噪声和天气表情更新，后者阻止空闲动作调度器和拖拽动画。

\subsection{MotionAgent --- 自主动作决策}

\texttt{MotionAgent} 在 V2.1 中引入，负责自主判断何时执行动作、执行什么动作，基于本地 LLM（Qwen2.5:0.5b）决策：

\begin{itemize}
    \item \textbf{4 级动作密度}：High / Medium / Low / Sleep，根据当前焦点进程动态调整
    \item \textbf{焦点进程抑制}：检测到全屏游戏或编辑器在前台时自动降低动作频率
    \item \textbf{情绪状态机}：\texttt{EmotionState} 维护基础情绪（neutral/happy/sad/angry/surprised），影响动作选择
    \item \textbf{空闲弧检测}：长时间无交互时触发更频繁的随机动作
    \item \textbf{本地决策}：通过 \texttt{LocalLLMClient} 调用 Ollama，延迟可控
\end{itemize}

% ═══════════════════════════════════════════
\section{闭环具身智能验证}
% ═══════════════════════════════════════════

\subsection{架构设计}

闭环具身智能验证系统通过"生成 → 执行 → 观察 → 评估 → 优化"的循环，
使 AI 的动作能力持续改进：

\[
\text{DeepSeek(生成)} \rightarrow \text{Live2D(执行)} \rightarrow \text{截图(捕获)} \rightarrow \text{GLM-4V(评估)} \rightarrow \text{记忆(学习)}
\]

\subsection{MotionVerifier --- 动作验证器}

三级测试集共 19 个测试用例：

\begin{itemize}
    \item \textbf{Level 1 (对照组)}：5 个应走硬编码模板的动作（挥手、点头、摇头、鞠躬、伸懒腰）
    \item \textbf{Level 2 (测试组)}：10 个应走 LLM 翻译路径的动作
          （害羞捂脸、叉腰、惊讶捂嘴、忧郁远望、俏皮眨眼、行礼、缩团、骄傲抬头、歪头思考、合十祈祷）
    \item \textbf{Level 3 (边界组)}：4 个边界输入（空字符串、乱码、复杂指令、物理不可能指令）
\end{itemize}

评估指标包括：参数合规率、对称配比率、关键帧数、复位帧检测。

\subsection{VisionMotionVerifier --- 视觉验证器}

对每个测试用例执行完整流程：

\begin{enumerate}
    \item 走 LLM 路径生成动作
    \item 播放动作至峰值帧
    \item 截取 RenderTexture 为 PNG
    \item 发送截图 + 评价提示给 GLM-4V
    \item 解析 GLM 返回的评分 (1--5) 和评语
\end{enumerate}

GLM-4V 评价标准：

\[
\text{Score} = 
\begin{cases}
5, & \text{非常像，动作精准传达意图} \\
4, & \text{比较像，主要特征符合} \\
3, & \text{一般，部分特征符合} \\
2, & \text{不太像，很多特征缺失} \\
1, & \text{完全不像}
\end{cases}
\]

\subsection{DualModelValidator --- 双模型评分}

V2.1 中采用 GLM-4V 作为主要评分模型（pass 阈值 3/5），Qwen-VL-Plus 作为只记录不影响的辅助评分：

\begin{itemize}
    \item \textbf{GLM-4V}：主评分器，返回 1--5 分 + 文本评语
    \item \textbf{Qwen-VL-Plus}：仅用于日志记录，不参与结果判断
\end{itemize}

\subsection{闭环学习 (MotionMemoryManager)}

视觉验证结果反馈到 \texttt{MotionMemoryManager}：

\begin{itemize}
    \item 评分 $\geq 3$（通过阈值）：加入运动记忆库（上限 30 条），供后续生成参考
    \item 评分 $< 3$：记录到 VERIFIED FEEDBACK 列表（上限 10 条），作为负面示例附带
    \item \textbf{无望检测}：同一动作 5 次尝试评分均 $\leq 2/5$ 时标记为「无望」，不再尝试
    \item 记忆注入优先级：高评分 > 低评分，新注入按分数排序
    \item 下次 DeepSeek 调用时附带上一次 FEEDBACK 作为负面示例，引导改进
\end{itemize}

% ═══════════════════════════════════════════
\section{交互系统}
% ═══════════════════════════════════════════

\subsection{拖拽系统 (DragHandler)}

\texttt{DragHandler} 管理所有交互与点击穿透的协调。每帧根据鼠标位置和面板状态动态设置
\texttt{WS\_EX\_TRANSPARENT} 或清除该标记：

\begin{itemize}
    \item 鼠标在 \texttt{BallPanel} 交互区域 → 关闭穿透（允许点击面板）
    \item 鼠标在 \texttt{RightPanel} 交互区域 → 关闭穿透（允许点击按钮/输入）
    \item 面板关闭时 → 强制重置穿透状态
    \item 拖拽释放时计算速度向量传递给 DesktopPet 物理引擎
\end{itemize}

拖拽释放时的抛掷物理：
\[
\mathbf{v}_{\text{throw}} = \frac{\sum_{i=1}^{n} (\mathbf{p}_{i} - \mathbf{p}_{i-1})}{n \cdot \Delta t},\quad
|\mathbf{v}_{\text{throw}}| \leq \text{maxThrowSpeed}=12
\]

抛掷速度经多帧缓冲平均后限幅，防止瞬间抖动的异常高速。

\subsection{悬浮球辐射菜单 (BallPanel)}

V2.1 全新引入的交互方式。桌面右下角粉色 ✦ 悬浮球，单击展开辐射菜单：

\begin{itemize}
    \item \textbf{⚙ 设置面板}：任务权重滑块调节 + 保存/清空，420×580px 浮动窗口
    \item \textbf{📊 报告面板}：MotionMemoryManager 演武心经学习报告展示
    \item \textbf{📋 便签面板}：待办/已完成管理，增删改查操作
\end{itemize}

面板通过 \texttt{OnGUI} 渲染，使用自定义 \texttt{GUIStyle}。标题栏可拖拽移动，右键关闭面板。
\texttt{PanelType} 枚举管理三种面板状态（None/Settings/Report/Reminders）。

\subsection{右侧 Widgets 面板 (RightPanel)}

Windows 11 Widgets 风格右侧滑出面板：

\begin{itemize}
    \item \texttt{\textasciitilde} 键切换展开/收起，鼠标划过右边缘自动展开
    \item 1 秒自动隐藏延迟（鼠标离开后计时）
    \item 展开宽度 220px，收起宽度 8px，滑动速度 10f/s
    \item 4 个工具按钮：聊（聊天聚焦）/ 设（打开设置）/ 签（打开便签）/ 告（打开报告）
    \item 底部集成输入栏
\end{itemize}

\texttt{IsPointInInteractiveArea(Vector2)} 供 DragHandler 判断点击穿透状态。

\subsection{分区点击反馈}

通过 \texttt{CubismRaycaster} 实现分区检测：

\begin{itemize}
    \item \textbf{Head} (高度 > 0.6)：摸头眯眼锁定动画
    \item \textbf{Body} (0.3~0.6)：捂胸惊讶
    \item \textbf{Leg} (< 0.3)：害羞踢腿
\end{itemize}

\subsection{拖拽挣扎动画}

拖拽过程中触发挣扎动画：双臂交替划水 + 双腿交替 + 身体扭动 + 慌张表情。
帧间速度经平滑滤波后直接写入 Cubism 参数，绕过物理系统的 0.8\,s 延迟，实现即时响应。

拖拽时 20+ 头发输出参数由帧间速度直接驱动，裙子（Param82/87/84）、
法盘（Param49/51/57/60）同步跟随，实现物理惯性拖尾效果。

% ═══════════════════════════════════════════
\section{感知与记忆系统}
% ═══════════════════════════════════════════

\subsection{法眼活动追踪 (ActivityTracker)}

每 2 秒通过 \texttt{GetForegroundWindow()} + \texttt{GetWindowText()} 获取前台窗口标题，
按关键词分类为编程、设计、游戏、学习、浏览、办公、通讯、音乐等。

通过 \texttt{EnumWindows} 扫描所有可见顶层窗口，构建多窗口环境摘要。
\texttt{BrowserTabReader} 通过 UI Automation API 读取主流浏览器标签页标题
（Chrome/Edge/Firefox/Opera/Brave/Vivaldi），无需安装插件。

\subsection{长期记忆系统 (PetMemory)}

\subsubsection{三层分层结构}

\begin{itemize}
    \item \textbf{核心事实}：用户偏好、重要承诺、关键信息（始终保留）
    \item \textbf{重要记忆 (Top-5)}：按重要性分数排序
    \item \textbf{近期琐事}：时间衰减，自动淘汰
\end{itemize}

记忆总量上限 30 条。核心事实永久保留，超出上限时优先淘汰低分琐事。

\subsubsection{记忆类别与重要性}

\begin{table}[H]
\centering
\caption{记忆类别与初始重要度}
\label{tab:memory}
\begin{tabular}{lll}
\toprule
\textbf{类别} & \textbf{说明} & \textbf{初始重要度} \\
\midrule
tool & 工具调用记录 & 0.3 \\
conversation & 对话摘要 & 0.5 \\
observation & 法眼观察 & 0.2 \\
reflection & 反思提炼 & 0.8 \\
\bottomrule
\end{tabular}
\end{table}

\subsubsection{反思反射}

累计重要性积分达阈值（默认 30）时触发 LLM 提炼高层次洞察：

\[
\sum \text{importance} \geq \text{THRESHOLD}=30 \Rightarrow \text{TriggerReflection()}
\]

同话题在冷却期间（默认 120\,s）不再重复记录。反思结果按 \texttt{reflection} 类别存入核心层，初始重要度 0.8。

\subsection{时间与天气系统 (TimeWeatherController)}

\subsubsection{双源天气获取}

QWeather API 优先（需 Key），失败时自动回退 wttr.in（无 Key 要求）。
解析天气类型：Clear / Cloudy / Overcast / Rain / Drizzle / Thunder / Snow / Fog。

\subsubsection{AI 天气语录}

当 QWeather 可用时，调用 DeepSeek 批量生成符玄风格的天气台词，
按天气类型缓存，随天气变化自动刷新。

% ═══════════════════════════════════════════
\section{桌面物理引擎}
% ═══════════════════════════════════════════

\subsection{核心物理模型}

简化的 2D 像素物理引擎：

\[
\begin{aligned}
v_y &= v_y + g \cdot \Delta t \\
y &= y + v_y \cdot \Delta t \\
\mathbf{v} &= \mathbf{v} \cdot \text{AIR\_RESISTANCE}
\end{aligned}
\]

地面碰撞：触地后速度反向乘以反弹系数，低于阈值时停止。

\subsection{地面任务状态机}

5 种地面行为通过加权随机切换：

\begin{table}[H]
\centering
\caption{地面任务状态机}
\label{tab:ground-task}
\begin{tabular}{lcc}
\toprule
\textbf{任务} & \textbf{权重(默认)} & \textbf{行为} \\
\midrule
MoveLeftEdge & 1 & 向左走到屏幕左边缘 \\
MoveRightEdge & 1 & 向右走到屏幕右边缘 \\
MoveLeftTime & 1 & 向左走随机时长 (2--4\,s) \\
MoveRightTime & 1 & 向右走随机时长 (2--4\,s) \\
StopTime & 6 & 停止随机时长 (6--15\,s) \\
\bottomrule
\end{tabular}
\end{table}

\subsection{多显示器支持}

通过 \texttt{VirtualScreen} 获取跨越所有显示器的完整桌面区域，
物理边界自动适配多显示器环境。

\subsection{睡眠唤醒检测}

通过帧间时间间隙检测系统休眠/唤醒：

\[
\Delta t_{\text{gap}} > \text{SLEEP\_DETECT\_THRESHOLD} \Rightarrow \text{RebuildDWM() + ResetPhysics()}
\]

% ═══════════════════════════════════════════
\section{窗口与系统集成}
% ═══════════════════════════════════════════

\subsection{Win32 API 封装}

\subsubsection{DWM 透明窗口}

窗口透明使用 \texttt{DwmExtendFrameIntoClientArea}，MARGINS 结构
\texttt{cxLeftWidth = -1} 表示全窗口玻璃层扩展。Camera 背景统一为纯黑 (0,0,0,0)。

\subsubsection{系统托盘}

使用 \texttt{Shell\_NotifyIcon} 实现系统托盘图标，支持左键切换显示/隐藏、
右键弹出菜单（显示/隐藏/开机自启/退出）。

\subsubsection{开机自启}

通过注册表实现：
\[
\text{HKCU}\backslash\text{Software}\backslash\text{Microsoft}\backslash\text{Windows}\backslash\text{CurrentVersion}\backslash\text{Run}
\]

\subsection{崩溃日志系统}

通过 \texttt{Application.logMessageReceived} 捕获异常和错误，追加到本地日志文件。
日志长度超限时自动截断（保留后半部分），防日志爆炸。

\subsection{进程互斥}

使用 Mutex（\texttt{DesktopPet\_Unity\_SingleInstance}）防止多实例。

% ═══════════════════════════════════════════
\section{性能优化}
% ═══════════════════════════════════════════

\subsection{三级性能档位}

\texttt{PerformanceMonitor} 管理三级 FPS 自适应档位：

\begin{table}[H]
\centering
\caption{性能档位}
\label{tab:perf}
\begin{tabular}{lccc}
\toprule
\textbf{档位} & \textbf{目标帧率} & \textbf{RT 缩放} & \textbf{触发条件} \\
\midrule
High & 60\,fps & 100\% & FPS > 45 且 CPU/GPU < 80\% \\
Normal & 40\,fps & 75\% & FPS 30--45 或 CPU/GPU < 85\% \\
Low & 20\,fps & 50\% & FPS < 30 或 CPU/GPU > 92\% \\
\bottomrule
\end{tabular}
\end{table}

\subsection{帧率采样与评估}

滚动 90 帧采样窗口，每 30 帧评估一次性能，决定是否升降档。

\subsection{CPU/GPU 监控}

\begin{itemize}
    \item \textbf{CPU}：通过 Win32 \texttt{GetSystemTimes} 获取总占用率
    \item \textbf{GPU}：通过 NVML（NVIDIA Management Library）获取 GPU 占用率
    \item \textbf{升档拦截}：CPU/GPU 占用 $> 80\%$ 时阻止升档
    \item \textbf{紧急降级}：CPU/GPU 占用 $> 92\%$ 时强制降至 Low 档
\end{itemize}

\subsection{内存管理}

\[
\begin{cases}
\text{appMemory} > 800\,\text{MB}, & \text{GC.Collect()} \\
\text{appMemory} > 1200\,\text{MB}, & \text{GC.Collect(2, Forced)} \\
\text{systemMemory} > 85\%, & \text{Resources.Unload + GC} \\
\text{systemMemory} > 93\%, & \text{强制 Low 档 + GC}
\end{cases}
\]

系统总内存通过 \texttt{GlobalMemoryStatusEx} 获取，应用内存通过 \texttt{Profiler.GetTotalAllocatedMemoryLong} 获取。

% ═══════════════════════════════════════════
\section{便签与提醒系统}
% ═══════════════════════════════════════════

\subsection{数据模型}

支持一次性、每日、工作日、每周四种重复规则。所有提醒序列化为 JSON 持久化到本地。

\subsection{推送链路}

到期提醒通过三级推送链路依次触发：
头顶气泡显示 $\rightarrow$ Windows Toast 通知 $\rightarrow$ Server酱³ 手机推送。

\subsection{去重机制}

AI 手动设置与服务器推送的同类提醒自动去重。创建前通过
\texttt{HasPendingReminderContaining()} 检查关键词，防止重复。

% ═══════════════════════════════════════════
\section{构建与部署}
% ═══════════════════════════════════════════

\subsection{构建脚本}

标准构建使用 \texttt{build.ps1}，支持 \texttt{-Quick} 参数仅验证编译（不输出可执行文件）：

\begin{lstlisting}[style=csharp,caption={BuildScript.cs 核心逻辑},label={lst:build}]
public static void BuildDesktopPet() {
    string[] scenes = EditorBuildSettings.scenes
        .Where(s => s.enabled).Select(s => s.path).ToArray();
    
    BuildPlayerOptions options = new BuildPlayerOptions {
        scenes = scenes,
        locationPathName = "Build/DesktopPet.exe",
        target = BuildTarget.StandaloneWindows64,
        options = BuildOptions.None
    };
    
    BuildPipeline.BuildPlayer(options);
}
\end{lstlisting}

\subsection{构建后验证}

\begin{itemize}
    \item 关键文件检查：\texttt{Build/DesktopPet.exe}、\texttt{Assembly-CSharp.dll}
    \item 构建日志关键字搜索：\texttt{error}、\texttt{warning}、\texttt{构建完成}
\end{itemize}

% ═══════════════════════════════════════════
\section{版本路线图}
% ═══════════════════════════════════════════

\begin{table}[H]
\centering
\caption{版本演进}
\label{tab:versions}
\begin{tabular}{llp{7cm}}
\toprule
\textbf{版本} & \textbf{日期} & \textbf{主要变更} \\
\midrule
v0.2 & --- & PNG 渲染初步 + API Key 环境变量 \\
v0.3 & --- & 10 种空闲动作 + 加权随机选择 \\
v0.4 & --- & 右键菜单 + 权重编辑 \\
v0.5 & --- & 动作锁定 + 走路淡入过渡 \\
v0.6 & 2026-06-13 & LaTeX 文档 + 右键菜单完善 \\
v0.7 & 2026-06-17 & 修复物理覆盖 + 执行顺序修正 \\
v0.8 & 2026-06-17 & 物理直接驱动 + 分区点击 + 挣扎动画 + 天气响应 \\
v0.9 & 2026-06-18 & DWM 玻璃层透明 + 底部输入栏 \\
N10 & 2026-06-18 & 点击穿透 + 右键菜单重构 + 多轮工具调用 \\
N11 & 2026-06-18 & 优先级气泡系统 + AI 回复稳定 \\
N12 & 2026-06-19 & 性能监控 + 开机自启 \\
N13 & 2026-06-20 & 服务端推送轮询 + Server酱³ 集成 \\
N14 & 2026-06-20 & 小程序课表数据打通 + 3 个学业查询工具 \\
N15 & 2026-06-20 & 修复逐句显示 bug \\
N16 & 2026-06-20 & Everything 毫秒级文件搜索 \\
N17 & 2026-06-22 & 提醒去重 + 服务器推送去重 \\
N18 & 2026-06-22 & 已完成任务独立视图 + 系统总内存监控 \\
N32d & 2026-07-07 & 闭环具身智能 + GLM-4V 视觉验证 + 运动记忆 \\
\textbf{V2.1} & \textbf{2026-07-09} & \textbf{悬浮球+辐射菜单 — BallPanel+RightPanel 取代右键菜单；
      动作系统重构到 Live2DFramework/ActionAgent/} \\
\bottomrule
\end{tabular}
\end{table}

% ═══════════════════════════════════════════
\section{附录}
% ═══════════════════════════════════════════

\subsection{环境变量清单}

\begin{table}[H]
\centering
\caption{环境变量}
\label{tab:env}
\begin{tabular}{lll}
\toprule
\textbf{变量名} & \textbf{必需} & \textbf{用途} \\
\midrule
DEEPSEEK\_API\_KEY & 是 & DeepSeek Chat API \\
GLM\_API\_KEY & 否 & GLM-4V 截图分析与动作验证 \\
QWEATHER\_API\_KEY & 否 & 和风天气数据源 \\
\bottomrule
\end{tabular}
\end{table}

\subsection{相关文件路径}

\begin{table}[H]
\centering
\caption{文件路径}
\label{tab:paths}
\begin{tabular}{ll}
\toprule
\textbf{用途} & \textbf{路径} \\
\midrule
构建日志 & \texttt{code/desktop\_unity/build\_log6.txt} \\
崩溃日志 & \texttt{Build/crash\_log.txt} \\
长期记忆 & \texttt{\{persistentDataPath\}/pet\_memory.json} \\
天机簿配置 & \texttt{\{persistentDataPath\}/pet\_config.json} \\
提醒数据 & \texttt{\{persistentDataPath\}/reminders.json} \\
系统提示词 & \texttt{Assets/Resources/SystemPrompt.txt} \\
空闲动作配置 & \texttt{Assets/Resources/IdleActions/idle\_actions.json} \\
浮游法阵图 & \texttt{Assets/Resources/magic\_circle.png} \\
\bottomrule
\end{tabular}
\end{table}

\subsection{关键设计决策}

\begin{table}[H]
\centering
\caption{设计决策}
\label{tab:decisions}
\begin{tabular}{lp{8cm}}
\toprule
\textbf{决策} & \textbf{理由} \\
\midrule
DWM 玻璃层而非色键 & 彻底解决半透明绿边问题 \\
DeepSeek Chat（非本地） & 高质量对话 + Function Calling 原生支持 \\
GLM-4V（非 GPT-4V） & 国内可用 + 图像理解质量达标 \\
LLM 翻译关键帧 & 无限动作多样性 vs 硬编码有限集 \\
JSON 配置驱动空闲动作 & 扩展性，无需改代码即可添加新动作 \\
三层分层记忆 & 核心稳定 + 重要不丢 + 琐事自动淘汰 \\
QWeather 优先 + wttr.in 回退 & 准确性 + 无需注册的双保险 \\
Everything CLI 优先 & 毫秒级 vs 分钟级的巨大差距 \\
悬浮球+右侧面板取代右键菜单 & 更好利用屏幕空间，右下方悬浮不遮挡桌面内容 \\
4 级密度动作决策 & 根据焦点进程抑制空闲动作，减少 AI 频繁介入干扰 \\
\bottomrule
\end{tabular}
\end{table}

\end{document}
